%=====================================================================
%  05_airloops_control.tex --- The scavenging air loop, the air-to-air
%  cooler, and the coupled control/solution procedure.
%
%  Code of record: geothermal_cooling/dbhe_master.ipynb (v10.6)
%    notebook cell 28 -> regen_air_pass        (scavenging pass + mixing node)
%    notebook cell 30 -> process_inlet_state   (wheel process inlet)
%    notebook cell 32 -> hatchery_loads, air_air_cooler
%    notebook cell 36 -> run_facility          (coupled driver, d_rA control)
%    notebook cell  6 -> configuration constants, Sections A--U
%  Baseline numbers: verification/reference/v10.6_baseline.json
%=====================================================================

\section{The scavenging air loop}
\label{sec:air:scav}

The wheel of Section~\ref{sec:dw:main} only moves moisture; something must
carry that moisture out of the building. That is the job of the scavenging
air stream: it is heated by the geothermal coil, passed through the wheel's
regeneration sector where it strips the sorbent, and then thrown away. The
word \emph{loop} survives from an earlier topology and is now slightly
generous, as Section~\ref{sec:air:fixedpoint} explains, but the stream is
still the one place where the whole plant's latent duty leaves the site.

\subsection{Stations and the code names that carry them}
\label{sec:air:stations}

One pass of the stream is the function \texttt{regen\_air\_pass}. Four states
matter, and the notebook names them as follows.

\begin{center}
\begin{tabular}{@{}llll@{}}
\toprule
Station & Symbol & Notebook name & Set by \\
\midrule
Scavenging inlet   & $(T_{\mathrm{ri}},w_{\mathrm{ri}})$ & \texttt{T\_ri, w\_ri} & the mixing node, downstream of the fan \\
Coil outlet        & $(T_{\mathrm{ro}},w_{\mathrm{ri}})$ & \texttt{T\_ro}        & the geothermal coil (sensible: $w$ unchanged) \\
Wheel exhaust      & $(T_{\mathrm{rd}},w_{\mathrm{rd}})$ & \texttt{T\_rd, w\_rd} & the wheel regeneration balance; \textbf{vented} \\
Mixed state        & $(T_{\mathrm{rm}},w_{\mathrm{rm}})$ & \texttt{T\_rm, w\_rm} & building exhaust $+$ outdoor air \\
\bottomrule
\end{tabular}
\end{center}

\noindent
Two further states enter from outside the stream: the outdoor air
$(T_{\mathrm{oa}},w_{\mathrm{oa}})$, and the building exhaust, which is the
room state $(T_{\mathrm{pc}},w_{\mathrm{pc}})$ because the building is held at
comfort and does not float.

\subsection{The pass, station by station}
\label{sec:air:pass}

\emph{Coil.} The scavenging air enters the geothermal coil at
$(T_{\mathrm{ri}},w_{\mathrm{ri}})$ and leaves at $T_{\mathrm{ro}}$, computed
by the $\varepsilon$--NTU relation of Eq.~\eqref{eq:hyd:coilout} with the
conductance $\UA_{\mathrm{coil}}$ of Eq.~\eqref{eq:hyd:UAcoil}. The coil is a
liquid-to-air exchanger with a dry surface, so it is purely sensible:
\begin{equation}
T_{\mathrm{ro}} = T_{\mathrm{ri}} + \frac{Q_{\mathrm{coil}}}{\mdr\,c_{p,a}},
\qquad
w_{\mathrm{ro}} = w_{\mathrm{ri}} .
\label{eq:air:coilpass}
\end{equation}
The code never forms $w_{\mathrm{ro}}$ at all; it passes \texttt{w\_ri}
straight into the wheel, which \emph{is} Eq.~\eqref{eq:air:coilpass} written
as an identity rather than an assignment.

\emph{Wheel.} The regeneration sector is entered at
$(T_{\mathrm{ro}},w_{\mathrm{ri}})$ against the process stream
$(T_{\mathrm{pi}},w_{\mathrm{pi}})$; the outlets follow from
Eqs.~\eqref{eq:dw:etadef}, \eqref{eq:dw:wrd} and~\eqref{eq:dw:hrd}. The wheel
delivers two things at once: the dried process air
$(T_{\mathrm{pd}},w_{\mathrm{pd}})$, which the building will use, and the wet
exhaust $(T_{\mathrm{rd}},w_{\mathrm{rd}})$, which the plant must lose.

\emph{Exhaust, mixing node, fan.} The wheel exhaust is vented in its entirety:
in v10.6 the pair $(T_{\mathrm{rd}},w_{\mathrm{rd}})$ is a \emph{reported
diagnostic and nothing else}, read by no downstream equation. That is the
single most important structural fact about this section, and
Section~\ref{sec:air:v105} explains why it was made so. A fresh scavenging
inlet is then assembled from two reservoirs --- building exhaust and outdoor
air --- and pushed through the fan.

\subsection{The mixing node}
\label{sec:air:mix}

Let $d_{r\mathrm{A}}$ be the fraction of the scavenging mass flow drawn from
the \emph{building exhaust}, and $d_{r\mathrm{B}}=1-d_{r\mathrm{A}}$ the
fraction drawn from \emph{outdoor air}. Applying the fractional mixing rule
of Eq.~\eqref{eq:psy:mix_frac} to the two conserved quantities --- water and
enthalpy --- and then recovering the temperature with
Eq.~\eqref{eq:psy:T_from_hw}:
\begin{align}
w_{\mathrm{rm}} &= d_{r\mathrm{A}}\,w_{\mathrm{pc}}
                 + d_{r\mathrm{B}}\,w_{\mathrm{oa}},
\label{eq:air:mix_w}\\[2pt]
h_{\mathrm{rm}} &= d_{r\mathrm{A}}\,h(T_{\mathrm{pc}},w_{\mathrm{pc}})
                 + d_{r\mathrm{B}}\,h(T_{\mathrm{oa}},w_{\mathrm{oa}}),
\label{eq:air:mix_h}\\[2pt]
T_{\mathrm{rm}} &= \frac{h_{\mathrm{rm}} - 2501.0\,w_{\mathrm{rm}}}
                        {1.006 + 1.86\,w_{\mathrm{rm}}} .
\label{eq:air:mix_T}
\end{align}
The fan is treated as isenthalpic, with the temperature-rise hook retained at
zero so that fan heat can be added later without renaming a state:
\begin{equation}
T_{\mathrm{ri}} = T_{\mathrm{rm}} + \Delta T_{\mathrm{fan}},
\qquad
w_{\mathrm{ri}} = w_{\mathrm{rm}},
\qquad
\Delta T_{\mathrm{fan}} = 0 .
\label{eq:air:fan}
\end{equation}

\paragraph{Both balances must use the same damper position.}
Equations~\eqref{eq:air:mix_w} and~\eqref{eq:air:mix_h} share one pair of
weights, and the code asserts it:
\verb|assert abs((d_rA_use + d_rB_use) - 1.0) < 1e-12|. The assertion looks
pedantic; it is not. A mixing junction produces a state that lies \emph{on the
straight line} joining its two inlet states on the psychrometric chart, and it
lies there precisely because water and enthalpy are apportioned with the same
two numbers. Use one damper position in the moisture balance and a different
one in the energy balance and the result is off the mixing line: a state that
no mixing process can produce, because mass and energy have been created in
different amounts.

An earlier version did exactly that, with weights that summed to $1.07$. The
error is not academic. Reconstructing it at the design point with weights
$(0.30,\,0.77)$ instead of $(0.30,\,0.70)$: the mixed humidity rises from
\SI{20.91}{\gram\per\kilogram} to \SI{22.41}{\gram\per\kilogram} and the mixed
enthalpy from \num{86.5} to \SI{92.6}{\kilo\joule\per\kilogram}, which lands
the scavenging inlet about \SI{2}{\kelvin} too hot and
\SI{1.5}{\gram\per\kilogram} too wet --- seven per cent of an air stream
conjured out of nothing, feeding a wheel whose whole output is a difference of
humidities. Errors of this kind do not announce themselves: the fixed point
still converges, the controller still reports success, and the plant appears
slightly better than it is.

\subsection{The physical cap on the exhaust fraction}
\label{sec:air:cap}

The building cannot exhaust more than it admits. Its fresh-air intake is the
forced ventilation $\mdv$ set by the eggs' oxygen requirement
(\SI{0.030}{\metre\cubed\per\hour} per egg, Section~S of the configuration),
and everything else in the process stream is recirculated and never leaves the
room. The exhaust available to the scavenging stream is therefore $\mdv$, and
the largest exhaust fraction physically obtainable is
\begin{equation}
\boxed{\;
d_{r\mathrm{A}}^{\mathrm{cap}} = \min\!\left(1,\ \frac{\mdv}{\mdr}\right)
\;}
\label{eq:air:cap}
\end{equation}
with $d_{r\mathrm{A}}^{\mathrm{use}} = \min(d_{r\mathrm{A}},
d_{r\mathrm{A}}^{\mathrm{cap}})$ and the balance made up with outdoor air. The
code applies the cap inside \texttt{regen\_air\_pass} and warns once per run
when the requested position exceeds it.

At the shipped configuration $\mdv=\SI{10.0}{\kilogram\per\second}$ and
$\mdr=0.70\,\mdp=\SI{23.52}{\kilogram\per\second}$, so
$d_{r\mathrm{A}}^{\mathrm{cap}}=\num{0.425}$. The nominal damper range
$[0.05,1]$ is thus a fiction above $0.425$: the scavenging stream is more than
twice the building's exhaust, and no damper setting can make the wheel breathe
more room air than the room produces. This has a direct consequence for the
controller (Section~\ref{sec:air:latent}) and it is worth stating as a design
observation in its own right: \emph{the authority of this actuator is set by
the ratio of two flows that were sized for entirely unrelated reasons} ---
egg respiration on one side, the wheel's regeneration ratio on the other.

\subsection{What changed in v10.5, and why it had to}
\label{sec:air:v105}

Until v10.4 the node of Section~\ref{sec:air:mix} mixed the \emph{wheel's own
exhaust} $(T_{\mathrm{rd}},w_{\mathrm{rd}})$ with outdoor makeup: a genuinely
recirculating loop, in which a fraction $d_{r\mathrm{A}}$ was vented and the
rest returned to the coil. That arrangement has a closed-form steady state
worth deriving, because it shows exactly how it fails.

Write the moisture the scavenging stream picks up in one pass, from
Eq.~\eqref{eq:dw:wrd},
\begin{equation}
\Delta w_{\mathrm{dw}} \equiv \frac{\mdp}{\mdr}\,
\left(w_{\mathrm{pi}}-w_{\mathrm{pd}}\right),
\qquad\text{so}\qquad
w_{\mathrm{rd}} = w_{\mathrm{ri}} + \Delta w_{\mathrm{dw}} .
\label{eq:air:dwdw}
\end{equation}
In the old topology the mixed humidity was
$w_{\mathrm{rm}} = d_{r\mathrm{B}}\,w_{\mathrm{rd}} + d_{r\mathrm{A}}\,
w_{\mathrm{oa}}$, and at steady state $w_{\mathrm{ri}}=w_{\mathrm{rm}}$, so
\begin{align*}
w_{\mathrm{ri}}
 &= d_{r\mathrm{B}}\left(w_{\mathrm{ri}}+\Delta w_{\mathrm{dw}}\right)
   + d_{r\mathrm{A}}\,w_{\mathrm{oa}}\\
w_{\mathrm{ri}}\left(1-d_{r\mathrm{B}}\right)
 &= d_{r\mathrm{B}}\,\Delta w_{\mathrm{dw}} + d_{r\mathrm{A}}\,w_{\mathrm{oa}},
\end{align*}
and since $1-d_{r\mathrm{B}}=d_{r\mathrm{A}}$,
\begin{equation}
\boxed{\;
w_{\mathrm{ri}} = w_{\mathrm{oa}}
+ \frac{d_{r\mathrm{B}}}{d_{r\mathrm{A}}}\,\Delta w_{\mathrm{dw}}
\;}
\label{eq:air:buildup}
\end{equation}

The moisture the wheel had just removed was being handed back to it, amplified
by the recirculation-to-vent ratio. Put the design numbers in. With
$\Delta w_{\mathrm{dw}} = (33.6/23.52)\times\SI{2.34}{\gram\per\kilogram}$ (the
year-1 value; it relaxes to \SI{1.85}{\gram\per\kilogram} by year 20)
$
= \SI{3.34}{\gram\per\kilogram}$ and the damper at its floor
$d_{r\mathrm{A}}=0.05$, Eq.~\eqref{eq:air:buildup} gives
$w_{\mathrm{ri}}=\SI{85}{\gram\per\kilogram}$ --- against a saturation humidity
of \SI{83}{\gram\per\kilogram} at the coil outlet temperature of
\SI{49.4}{\celsius}. The regeneration air arrived at the wheel supersaturated.

That is the failure, and it is a failure of the \emph{right} variable. A
desiccant wheel is not driven by humidity ratio but by the difference in
relative humidity between its two faces: the $F_2$ potential of
Eq.~\eqref{eq:dw:F2} is a relative-humidity-like coordinate, and the wheel
dries the process air only while the regeneration side sits at lower $\RH$
than the process side. The old loop drove that difference to zero and beyond.
The comparison at the design point is stark:

\begin{center}
\begin{tabular}{@{}lccc@{}}
\toprule
Regeneration inlet at $T_{\mathrm{ro}}=\SI{49.4}{\celsius}$
 & $w_{\mathrm{ri}}$ [\si{\gram\per\kilogram}] & $\RH$ & Wheel \\
\midrule
v10.6, exhaust $+$ outdoor mixing ($d_{r\mathrm{A}}=0.05$) & 21.3 & \SI{28}{\percent} & dries \\
Pre-v10.5, own-exhaust recirculation ($d_{r\mathrm{A}}=0.05$) & 84.9 & \SI{102}{\percent} & reverses \\
Pre-v10.5, own-exhaust recirculation ($d_{r\mathrm{A}}=0.30$) & 29.2 & \SI{38}{\percent} & dries weakly \\
\bottomrule
\end{tabular}
\end{center}

\noindent
The process face sits at $\RH\approx\SI{56}{\percent}$, so the current
topology opens a gap of some \num{28} percentage points, whereas the old one
closed it entirely at the very damper positions the controller wanted to use.
The wheel's regime diagnostic (Section~\ref{sec:dw:regime}) was written
precisely to catch the resulting reversal.

Since v10.5 the node draws from two reservoirs that the wheel does not
contaminate: room exhaust, which at comfort is drier in relative terms than
outdoor air, and outdoor air itself. The wheel's own product is vented and
never returns. The cost of the change is honest and should be stated: the plant
now heats air it has never heated before on every pass, so the coil duty is
higher than a recirculating loop would demand, and the exhaust enthalpy is
thrown away entirely. What is bought is a wheel that is not fighting itself.

\subsection{Does this stream still need a fixed point?}
\label{sec:air:fixedpoint}

The question is worth asking explicitly, because the code still solves
$(T_{\mathrm{ri}},w_{\mathrm{ri}})$ as unknowns. Trace the dependency in the
old topology:
$T_{\mathrm{ri}}\to T_{\mathrm{ro}}$ (coil)
$\to (T_{\mathrm{rd}},w_{\mathrm{rd}})$ (wheel)
$\to (T_{\mathrm{rm}},w_{\mathrm{rm}})$ (mix)
$\to (T_{\mathrm{ri}},w_{\mathrm{ri}})$ (fan). The chain closed on itself, the
inlet appeared on both sides, and a genuine two-variable fixed point had to be
solved.

Now trace it in v10.6. Equations~\eqref{eq:air:mix_w}--\eqref{eq:air:fan}
depend on $(T_{\mathrm{pc}},w_{\mathrm{pc}})$, on
$(T_{\mathrm{oa}},w_{\mathrm{oa}})$ and on $d_{r\mathrm{A}}$ --- all
prescribed --- and on \emph{nothing produced by the wheel or the coil}.
\textbf{Venting the exhaust severed the feedback}, in exactly the way the
building's comfort reset severs it on the process side. For a fixed
$d_{r\mathrm{A}}$ the scavenging inlet is a constant.

The baseline run confirms it: over twenty annual samples in
\texttt{v10.6\_baseline.json}, $T_{\mathrm{ri}}=\SI{31.299239}{\celsius}$ and
$w_{\mathrm{ri}}=\SI{21.289}{\gram\per\kilogram}$ to six figures at every
sample, while $T_{\mathrm{ret}}$, $T_{\mathrm{ro}}$, $T_{\mathrm{pd}}$ and
$\UA_{\mathrm{cool}}$ all drift with the depleting formation.

What survives is the \emph{water} fixed point of
Section~\ref{sec:hyd:circuit}: $T_{\mathrm{wi}}$ still closes on itself through
the well, the accumulator and the coil, and the scavenging inlet enters that
loop as the coil's air-side boundary condition. The driver carries
$(T_{\mathrm{ri}},w_{\mathrm{ri}})$ in the same iteration vector as
$T_{\mathrm{wi}}$, under-relaxed by one half,
\begin{equation}
x \leftarrow \tfrac12 x + \tfrac12 x^{\mathrm{new}},
\qquad
x \in \{T_{\mathrm{wi}},\,T_{\mathrm{ri}},\,w_{\mathrm{ri}}\},
\label{eq:air:relax}
\end{equation}
so their errors merely halve each sweep towards a target that never moves.
Keeping them in the vector costs almost nothing and preserves the machinery
for the day the topology changes again --- exhaust-air heat recovery on the
scavenging stream, for instance, would restore the coupling immediately.

\subsection{Verification: the mass balance}
\label{sec:air:massbal}

The node conserves dry air identically and by construction. The exhaust branch
carries $d_{r\mathrm{A}}^{\mathrm{use}}\mdr$, the outdoor branch
$d_{r\mathrm{B}}^{\mathrm{use}}\mdr$, and
\begin{equation}
d_{r\mathrm{A}}^{\mathrm{use}} + d_{r\mathrm{B}}^{\mathrm{use}} = 1
\quad\Longrightarrow\quad
\left(d_{r\mathrm{A}}^{\mathrm{use}}+d_{r\mathrm{B}}^{\mathrm{use}}\right)\mdr
= \mdr ,
\label{eq:air:massbal}
\end{equation}
which the code checks with the assertion quoted in
Section~\ref{sec:air:mix}. Two further balances close outside the node: the
wheel exhaust leaves at $\mdr$, matching what the node admits, and the
building exhaust drawn by the scavenging stream,
$d_{r\mathrm{A}}^{\mathrm{use}}\mdr \le \mdv$, is bounded by
Eq.~\eqref{eq:air:cap}. The assertion is a run-time check, not a printed
report; it fires as an exception rather than as a diagnostic line.

%=====================================================================
\section{The air-to-air cooler}
\label{sec:air:cooler}

\subsection{Role and configuration}
\label{sec:air:coolerrole}

The wheel returns the process air dry but hot: adsorption is exothermic, and
the latent load it removed reappears as sensible heat
(Eq.~\eqref{eq:dw:latent2sensible}). At the design point the wheel delivers
\SI{42.6}{\celsius} into a room held at \SI{37}{\celsius}. Something must take
that penalty back out before the air is supplied.

That something is \texttt{air\_air\_cooler}: an air-to-air, dry-surface,
$\varepsilon$--NTU exchanger. The hot side is the process stream
$(T_{\mathrm{pd}},w_{\mathrm{pd}})$ at $\mdp$; the cold side is a scavenging
stream drawn from outdoors at $T_{\mathrm{oa}}$ with flow
$\md_{\mathrm{cool}} = 1.5\,\mdp = \SI{50.4}{\kilogram\per\second}$
(\texttt{cool\_pro\_ratio}, Section~T of the configuration).

Two things follow, and they set everything else in this section. First, it is a
\emph{recuperator}: it moves heat from one air stream to another and consumes
no refrigeration, no water and no geothermal energy --- only fan power. It
cannot manufacture cold; it can only pass heat down a gradient it is given.
Second, being dry, it is sensible only:
\begin{equation}
\boxed{\;w_{\mathrm{ps}} = w_{\mathrm{pd}}\;}
\label{eq:air:cooler_w}
\end{equation}
where $(T_{\mathrm{ps}},w_{\mathrm{ps}})$ denotes the supply state delivered to
the building. (The notebook overloads the name \texttt{T\_pi}, using it both
for the wheel's process inlet and for the cooler outlet; the report keeps them
apart, and \texttt{T\_pi\_req}, \texttt{w\_pi\_req} in the driver are
$T_{\mathrm{ps}}^{\mathrm{req}}$, $w_{\mathrm{ps}}^{\mathrm{req}}$ here.)

Equation~\eqref{eq:air:cooler_w} is the hinge of the whole control system: the
last device before the building can change temperature and cannot change
humidity. Section~\ref{sec:air:principle} draws the consequence.

\subsection{Rating equations}
\label{sec:air:coolerrate}

With the moist-air specific heat of Eq.~\eqref{eq:psy:cp_moist} on the hot side
and dry-air $c_{p,a}$ on the cold side,
\begin{equation}
C_h = \mdp\left(1.006 + 1.86\,w_{\mathrm{pd}}\right)\times\num{1000},
\qquad
C_c = \md_{\mathrm{cool}}\,c_{p,a},
\qquad
C_r = \frac{C_{\min}}{C_{\max}} ,
\label{eq:air:caprates}
\end{equation}
and with $\varepsilon$ from the unmixed-crossflow relation
Eq.~\eqref{eq:hyd:epscross}, the duty and outlet are
\begin{equation}
Q_{\mathrm{cool}} = \varepsilon\,C_{\min}
\left(T_{\mathrm{pd}}-T_{\mathrm{oa}}\right),
\qquad
T_{\mathrm{ps}} = T_{\mathrm{pd}} - \frac{Q_{\mathrm{cool}}}{C_h} .
\label{eq:air:cooler_rate}
\end{equation}
At the design point $C_h=\SI{34.9}{\kilo\watt\per\kelvin}$ and
$C_c=\SI{50.7}{\kilo\watt\per\kelvin}$, so $C_{\min}=C_h$ and $C_r=0.689$: the
\emph{process} stream is the limiting one, which is why oversizing the ambient
fan buys progressively less.

The code also carries a trivial branch. If $T_{\mathrm{pd}}\le T_{\mathrm{oa}}$
it returns $T_{\mathrm{ps}}=T_{\mathrm{pd}}$ with $\UA=0$ and no saturation
flag: there is no gradient, so there is nothing to do and nothing to report.

\subsection{Target mode: back-solving the conductance}
\label{sec:air:target}

The cooler is not rated forward. The building balance, expressed through the
load interface as $T_{\mathrm{ps}}^{\mathrm{req}} = T_{\mathrm{pc}} -
\qs/(\mdp\,c_{p})$, states the supply temperature the space requires, and the
model asks what conductance delivers it. Inverting
Eq.~\eqref{eq:air:cooler_rate} for the effectiveness,
\begin{equation}
\boxed{\;
\varepsilon^{\mathrm{req}} =
\frac{C_h\left(T_{\mathrm{pd}}-T_{\mathrm{ps}}^{\mathrm{req}}\right)}
     {C_{\min}\left(T_{\mathrm{pd}}-T_{\mathrm{oa}}\right)}
\;}
\label{eq:air:eps_req}
\end{equation}
The crossflow relation cannot be inverted analytically for the NTU, so the code
bisects. Because $\varepsilon(\mathrm{NTU})$ increases monotonically, bisection
on $\mathrm{NTU}\in[0,\ \UA_{\max}/C_{\min}]$ is unconditionally convergent;
the code takes a fixed \num{60} halvings --- far more than needed, and cheap,
since each evaluation is one exponential --- and returns
\begin{equation}
\UA^{\mathrm{req}} = \mathrm{NTU}^{\mathrm{req}}\,C_{\min} .
\label{eq:air:UAreq}
\end{equation}
At the end of the first baseline year this gives
$\varepsilon^{\mathrm{req}}=\num{0.732}$, $\mathrm{NTU}=\num{2.42}$ and
$\UA_{\mathrm{cool}}=\SI{84.4}{\kilo\watt\per\kelvin}$, comfortably inside the
ceiling $\UA_{\max}=\SI{150}{\kilo\watt\per\kelvin}$.

Inverting the device this way is a modelling decision with a price. It reports
the conductance a real coil would need rather than the behaviour of a coil that
exists, so the model never tells you what happens at part load with a
\emph{fixed} exchanger --- the conductance simply follows the load. For sizing,
which is what this study is for, that is the more useful answer; for operation
it is not, and a fixed-$\UA$ rating mode should eventually sit beside it.

\subsection{The ambient floor}
\label{sec:air:floor}

A recuperator cannot cool its hot stream below the temperature of the stream it
rejects to. Heat flows down the gradient, and at $\varepsilon=1$ the outlet
reaches the cold inlet and stops:
\begin{equation}
\boxed{\;T_{\mathrm{ps}} \ge T_{\mathrm{oa}}\ \text{always}\;}
\label{eq:air:floor}
\end{equation}
The code enforces this by a sentinel rather than by a test: if
$T_{\mathrm{ps}}^{\mathrm{req}}\le T_{\mathrm{oa}}$ it sets
$\varepsilon^{\mathrm{req}}=2$, a value no exchanger can reach, which routes
execution to the saturated branch. Saturation therefore has two causes, and the
flag \texttt{sat\_cool} does not distinguish them:
\begin{enumerate}[nosep]
\item $T_{\mathrm{ps}}^{\mathrm{req}}\le T_{\mathrm{oa}}$ --- the target is
below ambient and \emph{no} dry recuperator can reach it, at any size;
\item $\varepsilon^{\mathrm{req}}>\varepsilon_{\max}$ --- the target is above
ambient but needs a coil larger than the ceiling.
\end{enumerate}
On either branch the block returns the best it can actually do at the ceiling,
\begin{equation}
T_{\mathrm{ps}} = T_{\mathrm{pd}}
- \frac{\varepsilon_{\max}\,C_{\min}\left(T_{\mathrm{pd}}-T_{\mathrm{oa}}\right)}
       {C_h},
\qquad \UA = \UA_{\max},
\label{eq:air:coolsat}
\end{equation}
rather than pretending the setpoint was met.

\paragraph{The design criterion.} Because the supply flow is chosen from the
design sensible load at a design depression $\Delta T_{\mathrm{supply}}$, the
required supply is $T_{\mathrm{ps}}^{\mathrm{req}} = T_{\mathrm{pc}} -
\Delta T_{\mathrm{supply}}$ by construction, and Eq.~\eqref{eq:air:floor}
reduces to a statement about two temperature differences:
\begin{equation}
\boxed{\;
\text{a dry ambient recuperator suffices}
\iff
\Delta T_{\mathrm{supply}} < T_{\mathrm{pc}} - T_{\mathrm{oa}}
\;}
\label{eq:air:coolercrit}
\end{equation}
At the shipped configuration $T_{\mathrm{pc}}-T_{\mathrm{oa}} =
37-31 = \SI{6}{\kelvin}$ and $\Delta T_{\mathrm{supply}}=\SI{3}{\kelvin}$, so
the criterion is met with a factor of two in hand and no below-ambient stage is
needed. The margin is a gift of the hatchery setpoint: a room held at
\SI{37}{\celsius} is warmer than the Gulf design day, which is what makes an
air-to-air recuperator viable at all. It is also fragile --- raise the ambient
design temperature to \SI{35}{\celsius}, or double the supply depression to
shrink the fan, and the criterion fails, at which point a below-ambient stage
(a dew-point evaporative cooler, or a chiller) becomes mandatory rather than
optional.

\paragraph{What the baseline actually shows.} \texttt{sat\_cool} does fire in
the opening days of a run, for a different reason: while the formation is
undepleted the coil delivers its hottest regeneration air, the wheel returns
$T_{\mathrm{pd}}\approx\SI{51}{\celsius}$, and cooling that to
\SI{34.1}{\celsius} against \SI{31}{\celsius} air needs more than
\SI{150}{\kilo\watt\per\kelvin}. In a one-year run the flag is raised for about
the first \num{10} steps (\num{2.5} days) and clears thereafter, with
$\UA_{\mathrm{cool}}$ falling from the ceiling to
\SI{98}{\kilo\watt\per\kelvin} within the first month. The binding constraint
in early life is the conductance ceiling, not the ambient floor --- and the
early-life excess of wheel duty is the same surplus that saturates the humidity
controller (Section~\ref{sec:air:latent}).

\subsection{Is there a condensation check? No}
\label{sec:air:condense}

There is none. Equation~\eqref{eq:air:cooler_w} is applied unconditionally:
\texttt{air\_air\_cooler} computes only temperatures, and the driver assigns
\verb|w_pi = ap['w_pd']| with no dew-point test anywhere in the function, in
the driver, or in the dashboard. A search of the notebook for a dew-point
variable returns nothing. If the block were ever driven below the stream's dew
point it would report a dry outlet at an impossible state, silently.

Whether that matters is a separate question, and here the answer is
reassuring. The supply carries $w_{\mathrm{pd}}=\SI{17.97}{\gram\per\kilogram}$
at the end of the first baseline year, a dew point of \SI{23.2}{\celsius},
while the outlet sits at \SI{34.1}{\celsius}: a margin of nearly
\SI{11}{\kelvin}. More strongly, Eq.~\eqref{eq:air:floor} guarantees
$T_{\mathrm{ps}}\ge T_{\mathrm{oa}}=\SI{31}{\celsius}$, so \emph{within this
topology} the cooler cannot reach the dew point of air the wheel has just
dried. The omission is safe by construction, not by luck.

It stops being safe the moment the cold side stops being ambient air. Add the
below-ambient stage that Eq.~\eqref{eq:air:coolercrit} would demand at a hotter
design day, or feed the cooler with chilled water from the absorption machine,
and the floor drops below the dew point; the block would then need a wet-coil
treatment (bypass factor or fin-efficiency), and $w_{\mathrm{ps}}<
w_{\mathrm{pd}}$ would become a second, uncontrolled dehumidification path.
This should be written before any below-ambient stage is added, not after.

%=====================================================================
\section{Control and the coupled solution}
\label{sec:air:control}

\subsection{The principle: control where there is an actuator}
\label{sec:air:principle}

The building imposes a two-component setpoint, $(T_{\mathrm{pc}},
w_{\mathrm{pc}})$, and the plant must satisfy two balances. Which of them can
be \emph{controlled} and which can only be \emph{checked} is not a modelling
preference; it is decided by where the hardware sits.

\begin{center}
\begin{tabular}{@{}lll@{}}
\toprule
Balance & Actuator downstream of the wheel? & Treatment \\
\midrule
Sensible & \textbf{Yes} --- the cooler & Control: back-solve $\UA_{\mathrm{cool}}$ \\
Latent   & \textbf{No} --- $w_{\mathrm{ps}}=w_{\mathrm{pd}}$, fixed at the wheel & Control from \emph{upstream}: $d_{r\mathrm{A}}$ \\
\bottomrule
\end{tabular}
\end{center}

Between the wheel outlet and the room there is exactly one device, and by
Eq.~\eqref{eq:air:cooler_w} it moves temperature only. Temperature is therefore
controlled locally, at the last possible moment; humidity has no local knob at
all and must be reached backwards, through the scavenging stream, the wheel and
in the last resort the well. The hierarchy is also an economic ordering ---
fan power first, a damper second, the geothermal resource last --- and that
ordering is the operational expression of what the project is trying to
conserve.

\subsection{The sensible channel}
\label{sec:air:sensible}

Nothing is left to do here beyond Section~\ref{sec:air:target}: the load
interface returns $T_{\mathrm{ps}}^{\mathrm{req}}$, the cooler is inverted for
$\UA^{\mathrm{req}}$, the floor test is applied and \texttt{sat\_cool} is set
if the target is unreachable. One scalar inversion, one bisection, no outer
iteration --- and, importantly, it happens \emph{after} the coupled solve,
because the cooler influences nothing upstream of itself.

\subsection{The latent channel}
\label{sec:air:latent}

The humidity setpoint is closed on the scavenging vent fraction. The model
solves
\begin{equation}
\text{find } d_{r\mathrm{A}} \in [d_{r\mathrm{A},\min},\,1]
\quad\text{such that}\quad
w_{\mathrm{pd}}\!\left(d_{r\mathrm{A}}\right)
= w_{\mathrm{ps}}^{\mathrm{req}} ,
\label{eq:air:rA_solve}
\end{equation}
with $d_{r\mathrm{A},\min}=\num{0.05}$ and a residual tolerance
$\num{1e-4}\,\si{\kilogram\per\kilogram}$ (\texttt{w\_ctrl\_tol}).

\paragraph{The mechanism.} Opening $d_{r\mathrm{A}}$ replaces outdoor air with
building exhaust in the scavenging stream. At the design point the room, held
at $\SI{37}{\celsius}$ and \SI{50}{\percent}, is drier than the outdoors both
absolutely (\num{19.9} against \SI{21.4}{\gram\per\kilogram}) and in relative
terms once heated, so more exhaust means drier regeneration air, a wider $F_2$
gap, and a drier process outlet: $\partial w_{\mathrm{pd}}/\partial
d_{r\mathrm{A}}<0$. Note carefully that this sign is a property of the
\emph{design weather}, not a theorem. On a cool dry morning the outdoor air is
the drier reservoir and the sign reverses. The code assumes the design sign in
its first probe step and would take one wasted iteration to recover if it were
wrong.

\paragraph{The scheme as coded.} It is a bounded secant, not the bisection its
predecessor used:
\begin{enumerate}[nosep]
\item Solve the coupled inner problem at the warm-started
      $d_{r\mathrm{A}}$; set $e_1 = w_{\mathrm{pd}}-w_{\mathrm{ps}}^{\mathrm{req}}$
      and seed $(r_0,e_0)=(d_{r\mathrm{A}},e_1)$.
\item Take the first probe as $r_1 = \mathrm{clip}\!\left(d_{r\mathrm{A}}
      \pm \num{0.1},\ d_{r\mathrm{A},\min},\ 1\right)$, the sign being
      $+$ if the air is too wet and $-$ if it is too dry.
\item Repeat at most $\num{12}$ times (\texttt{d\_rA\_ctrl\_max\_it}),
      stopping as soon as $|e_1| < \texttt{w\_ctrl\_tol}$: solve the inner
      problem at $r_1$, form $e_{\mathrm{new}}$, and take the secant step
      \begin{equation}
      r_{\mathrm{next}} = r_1 - e_{\mathrm{new}}\,
      \frac{r_1-r_0}{e_{\mathrm{new}}-e_0},
      \label{eq:air:secant}
      \end{equation}
      guarded by $|e_{\mathrm{new}}-e_0|>\num{1e-12}$ (below which the step is
      abandoned and $r_{\mathrm{next}}=r_1$), then shift
      $(r_0,e_0)\leftarrow(r_1,e_{\mathrm{new}})$ and
      $r_1\leftarrow\mathrm{clip}(r_{\mathrm{next}})$.
\item Adopt $d_{r\mathrm{A}}=r_0$ --- the last position actually
      \emph{evaluated}, never the untested extrapolation $r_1$ --- and re-solve
      the inner problem there so that every reported state is mutually
      consistent.
\end{enumerate}
Warm-starting from the previous time step matters: the loads move slowly, so
the previous solution is usually already inside tolerance.

\paragraph{Saturation at both ends.} Two flags mark the boundaries of the
actuator's authority:
\begin{align}
\texttt{sat\_hum} &:\quad d_{r\mathrm{A}} \ge 1-\num{1e-6}
   \ \ \text{and}\ \ w_{\mathrm{pd}} > w_{\mathrm{ps}}^{\mathrm{req}}
   + \texttt{w\_ctrl\_tol},
\label{eq:air:sathum}\\
\texttt{sat\_dry} &:\quad d_{r\mathrm{A}} \le d_{r\mathrm{A},\min}+\num{1e-6}
   \ \ \text{and}\ \ w_{\mathrm{pd}} < w_{\mathrm{ps}}^{\mathrm{req}}
   - \texttt{w\_ctrl\_tol}.
\label{eq:air:satdry}
\end{align}
The first is the classical failure: the damper is wide open, the air is still
too humid, the setpoint is unreachable at the current regeneration temperature,
and the only remaining escalation is to spend geothermal energy. The second was
added in v10.6 and is the failure that actually occurs: the damper is pinned at
its floor, the wheel is \emph{over}-drying, and the controller has no authority
left to stop it. Before v10.6 only \texttt{sat\_hum} existed, so v10.5 ran the
whole twenty years over-dry while the dashboard reported the humidity target
met. A one-sided flag on a two-sided actuator is not a diagnostic; it is a
blind spot.

Equation~\eqref{eq:air:sathum} carries a further weakness worth recording. The
controller's variable is bounded by $1$, but its \emph{effect} is bounded by
$d_{r\mathrm{A}}^{\mathrm{cap}}=\num{0.425}$ (Eq.~\eqref{eq:air:cap}). Between
\num{0.425} and \num{1} the residual is flat: the secant denominator collapses
below its guard, the update freezes, and the iterate only reaches exactly $1$
by accident. The wet-end flag can therefore be missed. The honest statement is
that the damper's authority ends at $\mdv/\mdr$, and the bound in
Eq.~\eqref{eq:air:rA_solve} should be $d_{r\mathrm{A}}^{\mathrm{cap}}$ rather
than $1$.

\subsection{What is coupled to what}
\label{sec:air:coupling}

\begin{center}
\begin{tabular}{@{}lll@{}}
\toprule
Sub-problem & Structure & Where it is solved \\
\midrule
Water loop            & closed on $T_{\mathrm{wi}}$      & inner fixed point \\
Scavenging inlet      & feed-through since v10.5 (Sec.~\ref{sec:air:fixedpoint}) & carried in the same vector \\
Process inlet         & feed-through (comfort reset)     & once, before the time loop \\
Latent control        & scalar root find on $d_{r\mathrm{A}}$ & outer loop \\
Sensible control      & scalar inversion for $\UA_{\mathrm{cool}}$ & after convergence \\
Formation, accumulator& stateful, marched                & once per step, at the end \\
\bottomrule
\end{tabular}
\end{center}

The water loop and the scavenging stream share the coil, so they are solved
\emph{simultaneously} in one iteration over
$(T_{\mathrm{wi}},T_{\mathrm{ri}},w_{\mathrm{ri}})$, never as nested loops ---
a full air solve inside every water iteration would multiply the cost for no
gain. Convergence is measured on
\begin{equation}
\max\left(\left|\Delta T_{\mathrm{wi}}\right|,\
          \left|\Delta T_{\mathrm{ri}}\right|,\
          \num{1e3}\left|\Delta w_{\mathrm{ri}}\right|\right)
< \texttt{fp\_tol} = \num{1e-4},
\label{eq:air:resid}
\end{equation}
with at most $\texttt{fp\_max}=\num{40}$ sweeps. The factor \num{1000} puts a
humidity ratio in \si{\kilogram\per\kilogram} on the same numerical scale as a
temperature in kelvin; without it the humidity residual would be met
automatically and would contribute nothing to the test.

\paragraph{Why the time loop lives in the driver.} Every block exposes
\texttt{init} / \texttt{set\_flow} / \texttt{eval} / \texttt{advance} and none
of them owns a clock. A feed-forward chain could be run block by block over the
whole horizon, but feedback and memory forbid it: the water loop's reinjection
temperature is an output of the same iteration that consumes it, and the
formation and the accumulator carry state that depends on every earlier step.
Both must be marched together, by something that sees all of them --- the
driver. The discipline that follows is strict and is respected in the code:
\texttt{dbhe\_advance} and \texttt{acc\_advance} are called \emph{once} per
step, with the converged state, after the controller has finished. Advancing
inside the iteration would commit history that the iteration was still free to
revise.

\subsection{Observed convergence}
\label{sec:air:converge}

Instrumenting the baseline case, one well evaluation per inner sweep:

\begin{center}
\begin{tabular}{@{}lr@{}}
\toprule
Quantity & Value \\
\midrule
Well evaluations, first step (cold start)      & \num{67} \\
Well evaluations, subsequent steps             & \num{24}--\num{25} \\
Inner solves per step (probe $+$ outer $+$ final) & \num{14} \\
Inner sweeps per solve, warm-started           & \num{1}--\num{2} \\
Recorded \texttt{fp\_it} (final consistent solve) & \num{1} \\
\bottomrule
\end{tabular}
\end{center}

The inner fixed point is well behaved: warm-started from the previous step, the
residual of Eq.~\eqref{eq:air:resid} is usually already met on the first sweep,
which is why the recorded \texttt{fp\_it} is $1$ throughout the baseline. That
number should be read for what it is --- the cost of the \emph{final} solve,
not the cost of the step.

The outer loop is the opposite story, and it is the diagnostic that matters. At
the shipped configuration the controller is pinned at $d_{r\mathrm{A},\min}$
with a residual of several \si{\gram\per\kilogram}, so the stopping test
$|e_1|<\texttt{w\_ctrl\_tol}$ is never met, the clipped secant returns the same
floor value every time, and the loop runs its full \num{12} iterations at every
one of the \num{29220} time steps. Roughly \num{12} of the \num{24} well
evaluations per step are spent proving, over and over, that a saturated
actuator is still saturated. A cheap fix exists --- break out when two
successive clipped iterates coincide --- and it would halve the run time
without changing a single result.

\subsection{Open item: the humidity actuator is the wrong one}
\label{sec:air:openitem}

\begin{center}
\fbox{\parbox{0.93\textwidth}{
\textbf{Planned for the next version: move humidity control from
$d_{r\mathrm{A}}$ to the coil-flow fraction.}\\[4pt]
The baseline run pins $d_{r\mathrm{A}}$ at its floor of \num{0.05} at every
step and raises \texttt{sat\_dry} throughout: the wheel removes
\SI{2.34}{\gram\per\kilogram} in year~1 and still
\SI{1.85}{\gram\per\kilogram} in year~20, while the building asks for
\SI{0.27}{\gram\per\kilogram}, and in the opening days it removes nearly
\SI{5}{\gram\per\kilogram}. The controller is saturated at its dry end from
the first step to the last, which means the plant is not being controlled at
all --- it is being run flat out and reported after the fact.\\[4pt]
The diagnosis is that $d_{r\mathrm{A}}$ is too weak an actuator for this duty.
It changes the \emph{quality} of the regeneration air, and by
Eq.~\eqref{eq:air:cap} it cannot change it by much, whereas the surplus to be
removed is a large fraction of the wheel's whole output. Meanwhile a large
share of the geothermal duty --- the coil heat that produces
$T_{\mathrm{ro}}$ --- is spent removing moisture the building does not want,
and the sensible penalty of that unwanted drying then has to be taken back out
by the cooler, at the cost of a bigger $\UA_{\mathrm{cool}}$.\\[4pt]
The next version therefore moves the latent control to the \emph{coil-flow
fraction}: throttling the geothermal water sent to the coil lowers
$T_{\mathrm{ro}}$ directly and therefore lowers the wheel's drying authority
continuously, from full duty down to none. It is a strong, monotone actuator
that acts on the resource itself, so under-drying and over-drying both become
reachable, and the geothermal energy not spent is geothermal energy saved ---
which, unlike a damper position, shows up as slower formation depletion. When
that change is made, $d_{r\mathrm{A}}$ should be retained as a secondary trim
and Eq.~\eqref{eq:air:rA_solve} rebounded on
$d_{r\mathrm{A}}^{\mathrm{cap}}$.}}
\end{center}
